\section{Introduction}

Semantic-ID (\sid) tokenization has become a practical interface for
generative recommendation. Systems such as TIGER map items into discrete code
sequences so that a sequence model can generate candidate identifiers rather
than score every catalog item directly~\cite{rajput2023tiger}. Once exported,
that item-to-code mapping becomes the address space exposed to the generator.
Recent work has expanded this space quickly: residual quantization frameworks,
recommender-native encoding, collaborative tokenization, differentiable
tokenization, non-uniform capacity, qualified collisions, variable-length
interfaces, and drift-aware refreshes all change the artifact that downstream
users must inspect before generator training~\cite{ju2025grid,liang2026resid,zhu2024cost,fu2026diger,wei2026card,hu2026quasid,cheng2026capsid,feng2026dact}.
This paper treats those exports as the object of inspection: the goal is not to
reimplement the literature, but to put item-to-code mappings under one
contract.

No comparable inspection layer exists for these exports. A tokenizer
repository may expose checkpoints, item mappings, codebooks, or only
intermediate feature files. Downstream users still need to know whether all
catalog items have valid codes, whether many items share a full code, whether
prefixes reflect user co-occurrence rather than only metadata taxonomy, whether
tail items receive address capacity, and whether the prefix structure creates
large fan-out. These checks are usually reimplemented ad hoc inside each paper,
making evidence hard to compare across codebook usage, downstream ranking, and
qualitative examples.

Aggregate ranking metrics remain necessary, but they do not expose the address
space itself. Prior recommender tooling has argued for behavioral tests beyond
NDCG-style aggregates~\cite{chia2022reclist}; \sid tokenizers add a concrete
resource problem: the item mapping should be validated and profiled before a
full generator consumes GPU time. Underused codebooks, repeated full codes,
behaviorally weak prefixes, tail compression, and large prefix fan-out can be
hidden by a downstream score, yet they are direct properties of the tokenizer
export.

We present \tool, a diagnostic/interface resource for inspecting item-to-\sid{}
tokenizer artifacts before or alongside downstream recommendation evaluation.
\tool is deliberately neither a new tokenizer nor a RecBole-style coverage
benchmark. It asks a narrower and reusable question: given any tokenizer
artifact that can emit item-level code sequences, what can be said about its
capacity use, aliasing, behavioral alignment, head-tail allocation, and
structural cost before retraining a full generator?

The contribution is an interface, a probe suite, and calibrated evidence. The
adapter contract covers \sid assignments, item metadata, interaction histories,
and optional generator outputs; validators reject incomplete or ambiguous rows
before metrics are reported. D1--D5 report utilization, aliasing, neighborhood
alignment, popularity allocation, and structural cost, with D6 churn support
for continual-tokenizer settings. Worked examples and controlled mechanism
probes show that full-code aliasing pressure, interaction-qualified aliasing
risk, collaborative-prefix recovery, tail capacity, and structural prefix
fan-out are separable. In the bounded Musical example, a category-prefix
control recovers more co-occurrence prefix neighbors than the learned export
rows, so learned addressability and behavioral prefix alignment require
separate inspection. Auxiliary vertical, fixed-reranker, and portability checks
support this finding without turning the paper into a method
leaderboard~\cite{cikm2026resource}.
